\documentclass{article}
\usepackage[utf8]{inputenc}
\usepackage[T1]{fontenc}
\usepackage{geometry}
\geometry{
    a4paper,
    left=2.5cm,
    right=2.5cm,
    top=2.5cm,
    bottom=2.5cm
}
\usepackage{graphicx}
\usepackage{float}
\usepackage{hyperref}
\usepackage[english]{babel}
\usepackage{amsmath}
\usepackage{amssymb}
\usepackage{array}
\usepackage{booktabs}
\usepackage{caption}
\usepackage{subcaption}
\usepackage{enumitem}
\usepackage{fancyhdr}
\usepackage{titlesec}
\usepackage{url}
\usepackage{textcomp}
\usepackage{xcolor}
\usepackage{listings}
\usepackage{multirow}
\usepackage{longtable}
\usepackage{tabularx}
\usepackage{times}

% Set up page style
\pagestyle{fancy}
\fancyhf{}
\rfoot{\thepage}
\renewcommand{\headrulewidth}{0pt}

% Title formatting
\titleformat{\chapter}[display]{\normalfont\bfseries}{\chaptertitle\thechapter}{0pt}{\Huge}
\titlespacing*{\chapter}{0pt}{50pt}{40pt}

% Document info
\title{Template Document}
\author{Academic Support Agent}
\date{\today}

\begin{document}
\maketitle

\maketitle

\begin{abstract}

%    \lipsum[1]



\par
{\centering
\includegraphics[width=1\linewidth]{ResumenEjemplo-MO-5to-Telec}
\par
}    


%\begin{figure}[H]
%    \centering
%    \includegraphics[width=0.5\linewidth]{figs/MiT_SaR.eps}
%    \caption{Caption}
%    \label{fig:my_label}
%\end{figure}
    

\end{abstract}





\section{Autocalificaci n}


\begin{itemize}
    \item Para su tarea borre del documento PDF todas las indicaciones.
    \item Antes de explicar las secciones de la presenta plantilla, se indica que para el presente documento se ha publicado una serie de v deos explicativos en:
    \begin{enumerate}
        \item 
    Cargar la Plantilla predise ada:
	 \href{https://youtu.be/FdkwJnJBmUg}{https://youtu.be/FdkwJnJBmUg}
        \item 
Edici n b sica la Plantilla:
	\href{https://youtu.be/xERlQ3jumR4}{https://youtu.be/xERlQ3jumR4}
        \item 
Insertar Im genes y modificar:
	\href{https://youtu.be/sq5S9gMj2xk}{https://youtu.be/sq5S9gMj2xk}
        \item 
Insertar referencias y modificar:
	\href{https://youtu.be/X3KJI-ExMYo}{https://youtu.be/X3KJI-ExMYo}
        \item 
Formato, Resaltada, otros:
	\href{https://youtu.be/TYsKYYPbkvU}{https://youtu.be/TYsKYYPbkvU}
    \end{enumerate}


\end{itemize}

\subsection{Indicaciones}
La Autocalificaci n consiste en una tabla que resume los Items o Ejercicios solicitados en la tarea. Esta tabla contiene una autocalificaci n de su tarea. 

Agregue una fila por cada  tem, actividad o ejercicio solicitado. Califique cada  tem con los siguientes porcentajes: 0\%, 25\%, 50\%, 75\% o 100\%.

\begin{table}[h]
    \caption{Ficha Resumen para autocalificaci n de la tarea}

    \centering

    \begin{tabular}{|c|c|}
        \hline

            \textbf{ tem, actividad o ejercicio solicitado} & \textbf{ Cumplimiento (0 a 100)}\\
        \hline
            Descripci n del Item 1 & 0 a 100 \\
        \hline
            Descripci n del Item 2 & 0 a 100 \\
        \hline
    \end{tabular}

    \label{tab:my_label}

\end{table}



%\lipsum[2-4]

%\begin{table}[h]
%    \caption{Etiqueta descriptiva de la Tabla}
%\centering
%    \begin{tabular}{|c|c|}
%        \hline
%            a & b \\
%        \hline
%            c & d \\
%        \hline
%    \end{tabular}
%    \label{tab:my_label}
%\end{table}

%Ver en la tabla \ref{tab:my_label}

%\lipsum[2-3]

\section{Desarrollo}


\subsection{Item 1: Escriba el  tem solicitado de la tarea}

Describa en computador (no a mano) como resolvi  lo solicitado en m nimo 2 l nea y m ximo 5 lineas. 

Enumere cada  tem, actividad o ejercicio solicitado y a continuaci n detalle la resoluci n. Si un  tem no realiz , incl yalo indicando porque no lo desarroll  (m ximo una l nea).

Agregue las im genes o capturas de pantalla de la tarea. En la Fig \ref{fig1_JoeCarrion} se muestra un ejemplo de imagen de captura de pantalla.

\begin{figure}[H]
    \centering
    \includegraphics[width=0.75\linewidth]{figs/Captura-Item1.png}
    \caption{Etiqueta descriptiva de la figura, este es un ejemplo de captura de pantalla de la tarea.}
    \label{fig1_JoeCarrion}
\end{figure}

\subsection{Item 2: Escriba el  tem solicitado de la tarea}

Describa en computador (no a mano) como resolvi  lo solicitado en m nimo 2 l nea y m ximo 5 lineas. 
Agregue las im genes o capturas de pantalla de la tarea. En la Fig \ref{fig2_JoeCarrion} se muestra un ejemplo de imagen de captura de pantalla. Mi texto

\begin{figure}[H]
    \centering
    \includegraphics[width=0.75\linewidth]{figs/Captura-Item2.png}
    \caption{Etiqueta descriptiva de la figura, reemplace con su propia descripci n.}
    \label{fig2_JoeCarrion}
\end{figure}

En la siguiente secci n se indica como utilizar referencias bibliogr ficas.

%test de cita \cite{katz1996handbook}

%Formato Equation , con numeracion


%\lipsum[1]
De acuerdo el autor \cite{koblas1992d}

\section{Indicaciones para referencias }

Siempre incluya las fuentes de consulta. En formato APA.
Puede utilice google Scholar como herramienta de ayuda. https://scholar.google.com/.


\subsection{Indicaciones para Citas}

Este es un ejemplo de una cita en Latex. De acuerdo al autor \cite{katz1996handbook} ...

Segundo ejemplo de una cita en Latex. La definici n es expuesta en \cite{koblas1992d} ...


Revise la siguiente imagen explicativa para uso de referencias bibliogr ficas. 
Utilice el formato indicado como APA en lugar del formato BitTex. 
Borre la imagen para entregar su tarea.


\begin{figure}[H]
    \centering
    \includegraphics[width=0.75\linewidth]{figs/Rubrica-Citas-Bittex.png}
    \caption{Indicaciones para inckuir referencias utilizando Google Scholar}
    \label{fig2_JoeCarrion}
\end{figure}

%






% Content from fichero1 would be included here






% Content from ficheroA would be included here
\end{document}